\section{Многослойный персептрон и обратное распространение ошибки}
\label{sec:mlp}

\noindent\hspace{1.25cm}\emph{Раздел используется в ЛР~№3.}

\subsection{Архитектура двухслойной сети}

Многослойный персептрон (multilayer perceptron, MLP)~--- сеть прямого распространения, в которой между
входом и выходным слоем расположены один или несколько \emph{скрытых} слоёв нелинейных нейронов. Для
аппроксимации функции одной переменной $y = f(t)$ достаточно сети 1--$S$--1: один вход, скрытый слой из
$S$ нейронов с функцией активации $\tanh$ и один линейный выходной нейрон (рисунок~\ref{fig:mlp}):
\begin{equation}
\label{eq:mlp}
\mathbf{a}_1 = \tanh(\mathbf{w}_1 t + \mathbf{b}_1), \qquad
y = \mathbf{w}_2^{\mathsf T}\mathbf{a}_1 + b_2 = \sum_{i=1}^{S} w_{2,i}\tanh(w_{1,i} t + b_{1,i}) + b_2 .
\end{equation}
Сеть имеет $3S + 1$ настраиваемых параметров. Каждый скрытый нейрон~--- это гладкая <<ступенька>>
$\tanh(w_{1,i}t + b_{1,i})$ с центром в точке $t = -b_{1,i}/w_{1,i}$ и крутизной $w_{1,i}$, а выходной
нейрон складывает ступеньки с весами $w_{2,i}$ (рисунок~\ref{fig:contributions}). Из суммы нескольких
ступенек можно собрать кривую любой формы; без нелинейности композиция линейных слоёв остаётся
линейной функцией, и сеть при любом $S$ даёт прямую.

\begin{figure}[!htb]
\centering
\begin{tikzpicture}[>=Stealth,
  neu/.style={circle, draw, minimum size=13mm, inner sep=0pt, font=\small},
  lbl/.style={font=\small, text=cgray, align=center}]
\node[neu, fill=cblue!12] (t) at (0, 0) {$t$};
\foreach \i/\y in {1/2.3, 2/0.8, S/-2.3} {
  \node[neu, fill=cyellow!25] (h\i) at (4.2, \y) {$\tanh$};
  \draw[->] (t) -- node[above, sloped, font=\footnotesize, pos=0.6] {$w_{1,\i}$} (h\i);
}
\node at (4.2, -0.65) {$\vdots$};
\node[neu, fill=cgreen!20] (o) at (8.4, 0) {$\Sigma$};
\foreach \i in {1, 2, S} \draw[->] (h\i) -- node[above, sloped, font=\footnotesize, pos=0.4] {$w_{2,\i}$} (o);
\draw[->] (o) -- ++(1.7, 0) node[right] {$y$};
\node[lbl] at (0, -3.35) {вход};
\node[lbl] at (4.2, -3.35) {скрытый слой, $S$ нейронов};
\node[lbl] at (8.4, -3.35) {выходной слой};
\end{tikzpicture}
\caption{Двухслойная сеть 1--$S$--1 (смещения не показаны)}
\label{fig:mlp}
\end{figure}

\emph{Теорема об универсальной аппроксимации} (Г.~Цыбенко, 1989~\cite{cybenko1989}; К.~Хорник, 1991)
утверждает: для любой непрерывной функции $f$ на компакте $K \subset \mathbb{R}^n$ и любого
$\varepsilon > 0$ существует сеть с \emph{одним} скрытым слоем из конечного числа сигмоидальных нейронов
и линейным выходом, для которой $\max_{\mathbf{x}\in K}|f(\mathbf{x}) - y(\mathbf{x})| < \varepsilon$.
Близкий по духу результат~--- теорема Колмогорова~--- Арнольда о представлении непрерывной функции
многих переменных суперпозицией функций одной переменной и сложения. Следствие для практики: двух слоёв
\emph{достаточно}, однако теорема не указывает нужное число нейронов, не гарантирует, что градиентный
спуск найдёт подходящие веса, и ничего не говорит о поведении сети между обучающими точками.

\begin{figure}[!htb]
\centering
\includegraphics[width=\textwidth]{\labfig{3}{10_neuron_contributions.png}}
\caption{Выход сети как сумма вкладов скрытых нейронов: функция
$f(t) = (t + 1,3)\sin(0,5t + 1)$ приближена тремя и пятью нейронами (ЛР~№3)}
\label{fig:contributions}
\end{figure}

\subsection{Алгоритм обратного распространения ошибки}

Обучение минимизирует MSE (\ref{eq:mse}) по всем весам и смещениям $\boldsymbol\theta$.
\emph{Обратное распространение ошибки} (backpropagation, Д.~Румельхарт, Дж.~Хинтон, Р.~Уильямс, 1986)~---
способ вычислить градиент $\nabla_{\boldsymbol\theta}E$ за один проход по сети от выхода к входу по правилу
дифференцирования сложной функции. Для сети (\ref{eq:mlp}) и выборки из $N$ точек:
\begin{enumerate}
\item \emph{прямой проход}: для каждого $t_k$ вычисляются $\mathbf{a}_{1,k}$ и выход $y_k$;
\item \emph{ошибка выхода}: $\delta_{2,k} = \frac{2}{N}(y_k - d_k)$, где $d_k = f(t_k)$~--- целевое значение;
\item \emph{обратный проход}: ошибка передаётся на скрытый слой через веса $\mathbf{w}_2$ и умножается на
производную активации, $\boldsymbol\delta_{1,k} = \delta_{2,k}\,\mathbf{w}_2 \odot (1 - \mathbf{a}_{1,k}^2)$
($\odot$~--- поэлементное произведение);
\item \emph{градиенты}: производная по весу равна сумме произведений сигнала ошибки на вход этого веса,
\begin{equation}
\begin{gathered}
\frac{\partial E}{\partial \mathbf{w}_2} = \sum_k \delta_{2,k}\mathbf{a}_{1,k}, \qquad
\frac{\partial E}{\partial b_2} = \sum_k \delta_{2,k}, \\
\frac{\partial E}{\partial \mathbf{w}_1} = \sum_k \boldsymbol\delta_{1,k} t_k, \qquad
\frac{\partial E}{\partial \mathbf{b}_1} = \sum_k \boldsymbol\delta_{1,k}.
\end{gathered}
\end{equation}
\end{enumerate}
Корректность реализации проверяют сравнением с численной производной
$\bigl(E(\theta + h) - E(\theta - h)\bigr)/2h$ при $h \sim 10^{-6}$: относительное расхождение должно
быть порядка $10^{-7}$ и меньше.

\subsection{Градиентные методы обучения}

\emph{Градиентный спуск} на каждой эпохе сдвигает параметры против градиента:
\begin{equation}
\label{eq:gd}
\boldsymbol\theta^{(k+1)} = \boldsymbol\theta^{(k)} - \eta\,\nabla E\bigl(\boldsymbol\theta^{(k)}\bigr),
\end{equation}
где $\eta$~--- коэффициент обучения (learning rate, \code{lr}). Вблизи минимума функцию ошибки можно приблизить
квадратичной формой с матрицей Гессе $H$. Тогда спуск устойчив только при
\begin{equation}
\label{eq:lr-bound}
\eta < \frac{2}{\lambda_{\max}},
\end{equation}
где $\lambda_{\max}$~--- наибольшее собственное число $H$, а скорость сходимости определяется
\emph{числом обусловленности} $\kappa = \lambda_{\max}/\lambda_{\min}$: чем больше $\kappa$, тем
сильнее поверхность ошибки вытянута в узкий <<овраг>>. Шаг приходится выбирать по самому крутому
направлению, и вдоль пологого дна спуск идёт очень медленно (рисунок~\ref{fig:gd-valley}). Для сетей
практикума $\kappa$ достигает $10^{11}$, поэтому градиентный спуск с постоянным шагом требует тысяч эпох.

\begin{figure}[!htb]
\centering
\includegraphics[width=\textwidth]{\figdir/gd_valley.png}
\caption{Градиентный спуск на вытянутой квадратичной поверхности ($\lambda_{\max} = 18$,
$2/\lambda_{\max} \approx 0,11$): 40 шагов из одной начальной точки}
\label{fig:gd-valley}
\end{figure}

Ускоренные алгоритмы (в скобках~--- реализации в Python):
\begin{itemize}
\item спуск с \emph{моментом} (\code{torch.optim.SGD(momentum=0.9)}):
$\Delta\boldsymbol\theta^{(k)} = \mu\,\Delta\boldsymbol\theta^{(k-1)} - \eta\nabla E$, $\mu \approx 0,9$. Инерция
гасит колебания поперёк оврага и накапливает скорость вдоль дна;
\item \emph{адаптивный шаг}: если ошибка уменьшилась, $\eta \leftarrow 1,05\,\eta$; если выросла более чем в
$1,04$ раза, шаг отменяется и $\eta \leftarrow 0,7\,\eta$ (правило \code{traingdx} MATLAB); близкий эффект
дают планировщик \code{torch.optim.lr\_scheduler.ReduceLROnPlateau} и \code{MLPRegressor(learning\_rate=}
\code{'adaptive')};
\item \emph{Adam} (\code{torch.optim.Adam}, \code{MLPRegressor(solver='adam')})~--- свой адаптивный шаг
для каждого параметра по скользящим средним градиента и его квадрата; стандарт для больших сетей;
\item метод \emph{Левенберга~--- Марквардта}~\cite{marquardt1963}, использующий якобиан $J$ вектора ошибок
сети по параметрам (\code{scipy.optimize.least\_squares(method='lm')}):
\begin{equation}
\label{eq:lm}
\Delta\boldsymbol\theta = -\bigl(J^{\mathsf T}J + \mu I\bigr)^{-1} J^{\mathsf T}\mathbf{e}.
\end{equation}
Матрица $J^{\mathsf T}J$ приближает гессиан, поэтому метод учитывает кривизну. При малом $\mu$ шаг
близок к шагу Гаусса~--- Ньютона, при большом~--- к короткому градиентному шагу; после удачного шага
$\mu$ уменьшается в 10 раз, после неудачного~--- увеличивается. Для сетей из десятков параметров метод
сходится за единицы--десятки эпох;
\item \emph{L-BFGS} (\code{MLPRegressor(solver='lbfgs')}, \code{torch.optim.LBFGS})~--- квазиньютоновский
метод, для малых сетей и выборок обычно самый быстрый.
\end{itemize}

\subsection{Предобработка, инициализация и критерии останова}

Входы и цели нормируют в диапазон $[-1; 1]$ (\code{MinMaxScaler} из \code{sklearn.preprocessing} с
параметром \code{feature\_range=(-1, 1)}):
\begin{equation}
\tilde t = 2\,\frac{t - t_{\min}}{t_{\max} - t_{\min}} - 1 ,
\end{equation}
чтобы аргументы $\tanh$ не попадали в область насыщения, где производная близка к нулю и обучение
останавливается. Выход сети пересчитывается обратно в исходный масштаб. Веса скрытого слоя
инициализируются методом Нгуена~--- Уидроу~\cite{nguyen1990}: активные области сигмоид равномерно
распределяются по отрезку входа, для одного входа $|w_{1,i}| = 0,7\,S$, а центры ступенек
$-b_{1,i}/w_{1,i}$ расставлены с равным шагом. В PyTorch слой \code{nn.Linear} по умолчанию берёт веса из
$U[-1/\sqrt{n}; 1/\sqrt{n}]$ ($n$~--- число входов), для $\tanh$ применяют инициализацию Глоро
(\code{nn.init.xavier\_uniform\_}).

Обучение прекращается при выполнении любого из условий:
\begin{itemize}
\item достигнуто максимальное число эпох \code{epochs} (в \code{MLPRegressor}~--- \code{max\_iter});
\item ошибка не превышает целевую: $E \le$ \code{goal};
\item норма градиента меньше порога (сеть <<застряла>> на плато);
\item ошибка на валидационной выборке не уменьшалась заданное число эпох~--- ранний останов
(\code{MLPRegressor(early\_stopping=True, validation\_fraction=0.15, n\_iter\_no\_change=6)}).
\end{itemize}
Параметр \code{show} задаёт период вывода протокола обучения и перерисовки графиков и на результат не
влияет.

\subsection{Обобщение: недообучение и переобучение}
\label{sec:generalization}

Число нейронов скрытого слоя $S$ определяет ёмкость сети. При малом $S$ сеть не может передать все изгибы
функции (недообучение): кривая слишком гладкая, экстремумы пропущены, ошибка велика и одинакова на
обучающих и контрольных точках. При избыточном $S$ на зашумлённых данных сеть проходит почти через
каждую обучающую точку, повторяя шум, и между точками возникают осцилляции и выбросы
(рисунок~\ref{fig:overfit}). Признак переобучения~--- ошибка на обучении ниже уровня шума при большой
ошибке на контрольных данных. В ЛР~№3 при шуме $\sigma = 0,7$ ошибка на обучающих точках с ростом $S$
монотонно убывает, а ошибка относительно истинной функции минимальна при $S = 3$ ($0,06$) и
возрастает до $4,06$ при $S = 30$; ранний останов снижает её до $0,37$.

\begin{figure}[!htb]
\centering
\includegraphics[width=\textwidth]{\labfig{3}{12_under_overfitting.png}}
\caption{Недообучение ($S = 2$), удачная сложность ($S = 3$) и переобучение ($S = 30$) на данных с
шумом $\sigma = 0,7$ (ЛР~№3)}
\label{fig:overfit}
\end{figure}

Переобучение следует отличать от \emph{недоученности} больших сетей при малом бюджете эпох: если ошибка
на контрольных точках равна ошибке на обучающих, рябь кривой исчезает при продолжении обучения.

\subsection{Реализация на Python}
\label{sec:py-mlp}

Одна эпоха пакетного градиентного спуска для сети (\ref{eq:mlp}) на нормированных данных \code{t},
\code{d} (массивы длины $N$) записывается векторно:
\begin{lstlisting}
a1 = np.tanh(t[:, None] * w1 + b1)        # N x S: выходы скрытого слоя
y = a1 @ w2 + b2                           # выход сети
d2 = 2 * (y - d) / len(t)                  # сигнал ошибки выходного слоя
d1 = np.outer(d2, w2) * (1 - a1 ** 2)      # ошибка, переданная на скрытый слой
w2 -= lr * a1.T @ d2;  b2 -= lr * d2.sum()
w1 -= lr * t @ d1;     b1 -= lr * d1.sum(axis=0)
\end{lstlisting}
Градиенты проверяют численно (\code{scipy.optimize.check\_grad}). Готовая сеть для сравнения~---
\code{MLPRegressor(hidden\_layer\_sizes=(10,), activation='tanh', solver='lbfgs')} или
\code{nn.Sequential(nn.Linear(1, S), nn.Tanh(), nn.Linear(S, 1))} в PyTorch, где градиенты
вычисляются автоматически (\code{loss.backward()})~\cite{pedregosa2011}.
